\section{Statistical Inference}
\label{sec:inference}

\subsection{Factorial Design Framework}

The factorial designs described in Section~\ref{sec:experiments} are analysed using ordinary least squares with effects coding ($x_i \in \{-1, +1\}$). This coding ensures orthogonality, meaning each main effect and two-way interaction can be estimated independently of all others, enabling clean attribution of outcome variation to specific factors. Experiment~1's $2^{10-1}$ half-fraction has Resolution~V, ensuring that all main effects and two-way interactions are estimable without bias from three-way terms, which are assumed negligible under effect sparsity.\footnote{In the smaller mixed-level designs (Experiments~2--3b), the degrees of freedom for error scale directly with replication count. Experiment~2's $3 \times 2^3$ design has 24 cells and 16 model parameters; with eight replicates the residual degrees of freedom are adequate. HC3 robust standard errors and wild bootstrap $p$-values (Appendix~\ref{sec:appendix_robustness}) confirm that findings are stable across all sub-experiments.} Replication in all experiments enables pure error estimation and lack-of-fit testing.

The stochastic variation in our experiments arises from seed-induced randomness in exploration, valuation draws, and tie-breaking. Inference is over the distribution of these random environments, not over a population of real markets. This is methodologically standard in computational economics and agent-based modeling \citep{Tesfatsion2006}. Factorial designs applied to stochastic simulations yield unbiased effect estimates under the same assumptions as physical experiments, with seed variation playing the role of experimental noise.

\subsection{Estimation}

For each response variable $Y$ (e.g., average revenue, convergence time, price volatility), we fit a linear model with main effects and all two-way interactions via ordinary least squares (OLS):
\begin{align}
  Y_i = \beta_0 + \sum_{j=1}^{k} \beta_j\, x_{ji} + \sum_{1 \le j < l \le k} \beta_{jl}\, x_{ji}\, x_{li} + \varepsilon_i,
  \label{eq:ols}
\end{align}
where $x_{ji} \in \{-1, +1\}$ is the coded level of factor $j$ in observation $i$, $\beta_j$ represents the main effect of factor $j$, and $\beta_{jl}$ captures the two-way interaction between factors $j$ and $l$. Under the coded parameterisation, $\beta_j$ equals half the change in mean response when factor $j$ moves from its low to its high level, holding all other factors at their centre values. We use Type~III ANOVA decomposition to assess the marginal contribution of each term,\footnote{Under the $\pm 1$ effects coding in a balanced factorial, the design matrix is orthogonal ($\mathbf{X}^\top\mathbf{X}$ is diagonal), so the OLS $t$-test for each coefficient is algebraically equivalent to the Type~III ANOVA $F$-test for the same term ($F = t^2$), yielding identical $p$-values. This equivalence depends on orthogonality; under $0/1$ dummy coding the columns would be correlated, main effect estimates would be conditional on the reference level rather than marginal, and the two tests would diverge.} testing each coefficient against the null $H_0\!: \beta = 0$ via the $t$-statistic $\hat{\beta}/\text{SE}(\hat{\beta})$.\footnote{The ANOVA $F$-tests are unadjusted for multiplicity. In the factorial design tradition \citep{Box2005}, the primary screening tools are the half-normal probability plot and Lenth's pseudo-standard-error method \citep{Lenth1989}, both of which identify active effects without requiring an independent error estimate. Formal multiplicity corrections (Holm--Bonferroni, Benjamini--Hochberg) are applied in the robustness analysis described below.} Model adequacy is validated by comparing OLS $R^2$ to a gradient-boosted machine (LightGBM) $R^2$ trained with five-fold cross-validation. Gaps below 0.05 support adequate linear approximation; gaps above 0.05 reveal detectable nonlinearity that warrants checking whether effect rankings change under the nonparametric model.

\subsection{Robustness and Model Adequacy}

We validate all findings through a comprehensive suite of robustness checks. For standard errors and multiple testing, we compute HC3 heteroskedasticity-consistent standard errors \citep{MacKinnonWhite1985}, Holm--Bonferroni and Benjamini--Hochberg corrections, and Rademacher wild bootstrap $p$-values for the top effects.

For specification and model adequacy, we run quantile regressions at five percentiles, lack-of-fit $F$-tests with pure error from replication, leave-one-out cross-validated $R^2$ via the PRESS statistic, LightGBM nonparametric benchmarking, and LASSO variable selection.

No multiplicity correction is applied across the six sub-experiments. The probability that the number of bidders ranks first among $k$ factors in all six independent sub-experiments by chance is bounded above by $(1/k)^6$, which for $k \geq 3$ is below 0.0014. This permutation argument provides conservative evidence that the ordinal ranking is unlikely to arise by chance.\footnote{$t$-statistic magnitudes are not comparable across experiments due to differences in sample size, model specification, and error variance; the comparison is strictly ordinal.}

We compute minimum detectable effects (MDE) at 80\% power and $\alpha = 0.05$:
\begin{align}
  \text{MDE}_{j} = \bigl(t_{\alpha/2,\,\nu} + z_{0.80}\bigr) \times \frac{\hat{\sigma}}{\sqrt{\sum_{i} x_{ij}^{2}}},
  \label{eq:mde}
\end{align}
where $\hat{\sigma}$ is the residual standard deviation, $\nu$ is the residual degrees of freedom, and $\sum_{i} x_{ij}^{2}$ is the sum of squared coded values for factor~$j$. For binary $\pm 1$ factors this reduces to $\hat{\sigma}/\sqrt{n}$; for orthogonal polynomial contrasts the sum of squares differs by column (Appendix~\ref{sec:appendix_inference}). MDEs range from 2--5\% of the mean response, indicating that the designs have sufficient power to detect practically meaningful effects. Full methodological details appear in Appendix~\ref{sec:appendix_inference}, and per-experiment diagnostics are reported in Appendix~\ref{sec:appendix_robustness}.

\subsection{Global Sensitivity Analysis}

As a complement to the factorial ANOVA, we decompose output variance through the Sobol--Hoeffding framework \citep{Sobol2001}. The first-order Sobol' index $S_i$ measures the fraction of total output variance attributable to factor $i$ alone, while the total-order index $S_{T_i}$ captures both its direct contribution and all interaction effects involving that factor. The gap $S_{T_i} - S_i$ quantifies how much of a factor's influence operates through interactions rather than through its main effect, providing a concise measure of factor interdependence that the standard ANOVA table does not directly report.

For balanced designs with effects coding ($\pm 1$), the columns of the design matrix are mutually orthogonal, and the Type~III ANOVA sum of squares decomposition coincides exactly with the analytical Sobol--Hoeffding decomposition \citep{Saltelli2008}. The first-order index reduces to $S_i = \mathrm{SS}_i / \mathrm{SS}_{\mathrm{total}}$, giving exact closed-form Sobol' indices without Monte Carlo sampling. This equivalence means that the ANOVA-based effect rankings and the variance-based sensitivity rankings are algebraically identical for these designs, providing a dual interpretation of the same decomposition.

To validate the factor rankings beyond this analytical equivalence, we apply seven independent sensitivity methods to each response variable in each experiment. These include random forest permutation importance, SHAP values via TreeSHAP, Morris elementary effects \citep{Morris1991}, Monte Carlo Sobol' indices computed on three distinct surrogate models, FAST, and Borgonovo's $\delta$ moment-independent measure \citep{Borgonovo2007}. Cross-method concordance is assessed by the mean pairwise Spearman rank correlation $\rho$ across all seven methods. High concordance ($\rho > 0.70$) indicates that factor rankings are robust to methodological assumptions; low concordance signals a flat importance landscape where minor factors are difficult to rank. Per-experiment Sobol' tables and concordance metrics are reported in Appendix~\ref{sec:appendix_sensitivity}.
